\documentclass[10pt,a4paper]{amsart}
\pdfinfoomitdate=1
\pdftrailerid{}
\pdfsuppressptexinfo=15

\usepackage[T1]{fontenc}
\usepackage{lmodern}
\usepackage[expansion=false]{microtype}
\usepackage[margin=25mm]{geometry}
\usepackage{amsmath,amssymb,mathtools}
\usepackage{xcolor}
\usepackage{booktabs}
\usepackage{array}
\usepackage{float}
\usepackage{enumitem}
\usepackage[numbers,sort&compress]{natbib}
\usepackage{hyperref}
\usepackage[nameinlink,noabbrev]{cleveref}

\hypersetup{
  colorlinks=true,
  linkcolor=blue!45!black,
  citecolor=blue!45!black,
  urlcolor=blue!45!black,
  pdftitle={},
  pdfauthor={},
  pdfsubject={},
  pdfkeywords={}
}

\newtheorem{theorem}{Theorem}
\newtheorem{lemma}[theorem]{Lemma}
\newtheorem{proposition}[theorem]{Proposition}
\newtheorem{corollary}[theorem]{Corollary}
\theoremstyle{remark}
\newtheorem{remark}[theorem]{Remark}

\newcommand{\Sub}{\operatorname{Sub}}
\newcommand{\im}{\operatorname{im}}
\newcommand{\Fix}{\operatorname{Fix}}
\newcommand{\depth}{\operatorname{depth}}
\newcolumntype{L}[1]{>{\raggedright\arraybackslash}p{#1}}

\title[Dihedral normalizer dynamics]{Iterated subgroup-normalizer dynamics\\
in finite dihedral groups}
\author{Anonymous}
\date{}
\subjclass[2020]{20D30, 37B15, 05C20, 11A25}
\keywords{dihedral group, subgroup normalizer, functional graph,
target fibre, arithmetic graph invariant}

\begin{document}

\begin{abstract}
Let \(D_{2n}=\langle r,s\mid r^n=s^2=1,\ srs=r^{-1}\rangle\), and
iterate \(H\mapsto N_{D_{2n}}(H)\) on the set of all subgroups.  The subgroup
coordinates and the one-step odd/even normalizer rule are classical and are
used here with full credit to their sources.  Writing \(n=2^a m\) with
\(m\) odd, we show that the remaining dynamics consists of
\(\sigma(m)\) full binary inverse trees of height \(a\), together with all
\(\tau(n)\) rotation subgroups feeding one distinguished fixed root.  We
derive every positive-time image and every target fibre.  We then prove that
the unlabelled functional graph is classified exactly by
\((a,\sigma(m),\tau(n))\).  In particular, the parameters \(33\) and \(35\)
give conjugate graphs, and so do all common power-of-two lifts.  Symbolic
proofs are paired with deterministic literal-normalizer replay.
\end{abstract}

\maketitle

\section{Scope, owned inputs, and theorem}

The subgroup structure of finite dihedral groups is classical
\citep{Cavior1975,ConradDihedralII}.  The normalizer of a
reflection-containing dihedral subgroup is likewise governed by a direct
odd/even rule; it appears in \citet{Frenkel2003} and in complete subgroup
cases in \citet{ShelashEtAl2023}.  Those results settle the carrier and the
one-step map.  They receive no contribution credit below.  Our subject is
what repeated application does to \emph{every} subgroup, including its full
inverse forest, all-time fibres, and the information about \(n\) retained by
the resulting unlabelled graph.

Put
\[
G_n=D_{2n}=\langle r,s\mid r^n=s^2=1,\ srs=r^{-1}\rangle,\qquad n\geq3.
\]
For \(d\mid n\) and \(0\leq j<d\), write
\[
R_d=\langle r^d\rangle,\qquad
H_{d,j}=\langle r^d,r^j s\rangle.
\]
The complete subgroup carrier is
\[
\Sub(G_n)=\{R_d:d\mid n\}\ \sqcup\
\{H_{d,j}:d\mid n,\ 0\leq j<d\},                       \tag{1}\label{eq:carrier}
\]
of size \(\tau(n)+\sigma(n)\).  Let
\(\mathcal N_n(H)=N_{G_n}(H)\), and write
\(n=2^a m\) with \(m\) odd.

\begin{theorem}[Full iterated normalizer graph]\label{thm:main}
For every \(n\geq3\), the map \(\mathcal N_n\) on \(\Sub(G_n)\) has the
following properties.
\begin{enumerate}[label=\textup{(\roman*)},leftmargin=*,itemsep=2pt]
\item The owned one-step rule and its iterates are
\[
\mathcal N_n(R_d)=H_{1,0},\qquad
\mathcal N_n(H_{d,j})
=H_{d/\gcd(d,2),\,j\bmod d/\gcd(d,2)},                 \tag{2}\label{eq:onestep}
\]
\[
\mathcal N_n^t(H_{2^k e,j})
=H_{2^{\max(k-t,0)}e,\,
     j\bmod 2^{\max(k-t,0)}e},                         \tag{3}\label{eq:iterate}
\]
for \(t\geq0\), \(0\leq k\leq a\), \(e\mid m\), and
\(0\leq j<2^ke\).
\item The dihedral states form \(\sigma(m)\) full binary inverse trees of
height \(a\), rooted at the fixed states \(H_{e,j}\) with
\(e\mid m\), \(0\leq j<e\).  Every \(R_d\) is an extra source of depth one
entering only \(H_{1,0}=G_n\).  Hence
\[
\depth(R_d)=1,\qquad \depth(H_{d,j})=v_2(d),
\]
and the depth polynomial is
\[
\mathcal D_n(z)=\sigma(m)+\tau(n)z+
\sigma(m)\sum_{k=1}^{a}2^kz^k.                         \tag{4}\label{eq:depth}
\]
All cycles are fixed points, and
\(\#\Fix(\mathcal N_n^t)=\sigma(m)\) for \(t\geq1\).
\item For every \(t\geq1\),
\[
|\im\mathcal N_n^t|
=\sigma(m)\bigl(2^{\max(a-t,0)+1}-1\bigr).             \tag{5}\label{eq:image}
\]
No rotation target has a positive-time preimage.  For every \(t\geq1\) and
every positive-level carrier target \(H_{2^ke,j}\) with
\(1\leq k\leq a\), \(e\mid m\), and \(0\leq j<2^ke\),
\[
\#(\mathcal N_n^t)^{-1}(H_{2^ke,j})
=
\begin{cases}
2^t,&k+t\leq a,\\
0,&k+t>a,
\end{cases}                                            \tag{6}\label{eq:fibrelevel}
\]
while a fixed root has fibre
\[
\#(\mathcal N_n^t)^{-1}(H_{e,j})
=2^{\min(t,a)+1}-1+\tau(n)\,\mathbf1_{(e,j)=(1,0)}.    \tag{7}\label{eq:fibreroot}
\]
\item Two such unlabelled directed functional graphs are isomorphic exactly
when
\[
\bigl(v_2(n),\sigma(\operatorname{odd}(n)),\tau(n)\bigr)
=\bigl(v_2(q),\sigma(\operatorname{odd}(q)),\tau(q)\bigr). \tag{8}\label{eq:signature}
\]
In particular, the graphs for \(n=33\) and \(n=35\) are conjugate, although
\(|G_{33}|\neq|G_{35}|\); the same holds for
\((2^b33,2^b35)\) for every \(b\geq0\).
\end{enumerate}
\end{theorem}

\paragraph{Proof dependency graph.}
The theorem is organized so the owner-covered bridge is visibly separated
from the dynamical deductions:
\[
\begin{gathered}
\boxed{\text{subgroup coordinates + one-step parity rule}}\\[-1mm]
\Downarrow\\[-1mm]
\boxed{\text{\(t\)-fold halving \eqref{eq:iterate}}}\\[-1mm]
\Downarrow\\[-1mm]
\text{binary forest}\Longrightarrow
\{\mathcal D_n,\im\mathcal N_n^t,\text{ target fibres}\},\\
\text{forest observables}\Longrightarrow\{a,\sigma(m),\tau(n)\}
\Longrightarrow\text{iff signature and collisions}.
\end{gathered}                                         \tag{9}\label{eq:dependency}
\]
Thus the signature theorem is not obtained merely from equal depth
polynomials; it uses the full component structure and total vertex count.

\section{From the one-step rule to the entire forest}

For completeness, we recall why \eqref{eq:carrier} and \eqref{eq:onestep}
hold.  A subgroup contained in \(\langle r\rangle\) is a unique \(R_d\).
Otherwise its rotation part is \(R_d\), and any reflection \(r^js\) in it
generates the unique nonrotation coset, giving \(H_{d,j}\) with \(j\bmod d\).
This also proves the count \(\tau(n)+\sigma(n)\).

Every \(R_d\) is normal in \(G_n\).  Conjugation by \(r^u\) sends
\(r^js\) to \(r^{j+2u}s\), so it normalizes \(H_{d,j}\) exactly when
\(d\mid2u\).  A general normalizing reflection \(r^us\) satisfies the
corresponding condition \(d\mid2(u-j)\).  Therefore all normalizers together
form
\[
\left\langle r^{d/\gcd(d,2)},r^js\right\rangle,
\]
which is precisely \eqref{eq:onestep}.  This derivation makes the coordinate
convention explicit, but the result itself is part of the credited input.

Write \(d=2^ke\) with \(e\) odd.  Iterating \eqref{eq:onestep} removes one
factor of two per step until none remains, and reduces \(j\) modulo the new
step.  This proves \eqref{eq:iterate}.  For a fixed pair
\((e,j_0)\), \(e\mid m\) and \(0\leq j_0<e\), the states at level \(k\) are
\[
H_{2^ke,j_0+\ell e},\qquad0\leq\ell<2^k.              \tag{10}\label{eq:lifts}
\]
Each non-leaf state has exactly two lifts at the next level.  Summing the
number of roots over odd divisors gives
\(\sum_{e\mid m}e=\sigma(m)\), so these are
\(\sigma(m)\) full binary trees of height \(a\).  Separately, all
\(\tau(n)\) rotations map directly to \(H_{1,0}\).  The claimed depths and
\eqref{eq:depth} follow.  Every orbit reaches a root, hence every recurrent
state is fixed.

\section{All-time images and target fibres}

At time \(t\geq1\), a level-\(\ell\) dihedral state can be reached exactly
when \(0\leq\ell\leq\max(a-t,0)\).  Each of the \(\sigma(m)\) trees then
contributes
\[
1+2+\cdots+2^{\max(a-t,0)}
=2^{\max(a-t,0)+1}-1
\]
image vertices, proving \eqref{eq:image}.  Rotations have no preimages
because every normalizer contains the subgroup it normalizes, whereas a
proper rotation target omits all reflections; this is also immediate from
\eqref{eq:onestep}.

Let \(H_{2^ke,j}\) be a carrier target with \(1\leq k\leq a\), \(e\mid m\),
and \(0\leq j<2^ke\).  A time-\(t\) source, for \(t\geq1\), must
occur at level \(k+t\).  When that level exists, the residue \(j\) has exactly
\(2^t\) lifts modulo \(2^{k+t}e\); otherwise there is no source.  This proves
\eqref{eq:fibrelevel}.  For a root, all levels \(0,\ldots,\min(t,a)\) have
already arrived, giving the geometric sum
\(2^{\min(t,a)+1}-1\).  In addition, every rotation has arrived at every
positive time, but only at the distinguished root.  This gives
\eqref{eq:fibreroot}.

These formulas also close source mass without appeal to graph intuition.
Indeed, summing all positive-level fibres, all root-tree fibres, and the
\(\tau(n)\) distinguished additions gives
\[
\sigma(n)+\tau(n)=|\Sub(G_n)|.
\]
Targets absent from \eqref{eq:image} have fibre zero.  Thus
\eqref{eq:fibrelevel}--\eqref{eq:fibreroot} are an every-target atlas rather
than only a distribution of indegrees.

As a concrete ledger, take \(n=12=2^2\cdot3\).  Here
\(a=2\), \(\sigma(m)=4\), \(\tau(n)=6\), and

\[
|\Sub(G_{12})|=\sigma(12)+\tau(12)=34,\qquad
\mathcal D_{12}(z)=4+14z+16z^2.                       \tag{11}\label{eq:n12depth}
\]
At time one the image has \(4(2^2-1)=12\) states.  Each of the eight
level-one targets has two sources; each ordinary root has three sources; and
the distinguished root has \(3+6=9\).  Thus

\[
8\cdot2+3\cdot3+9=34.
\]
At time two only the four roots remain.  The three ordinary roots have seven
sources each and the distinguished root has \(7+6=13\), giving
\(3\cdot7+13=34\).  This box displays both mechanisms in
\eqref{eq:fibreroot}: binary-tree mass grows geometrically, while the
rotation contribution is already saturated at time one.

Because every cycle is a fixed root, the finite-map zeta function follows as
\[
\zeta_{\mathcal N_n}(z)
=\exp\left(\sum_{t\geq1}\frac{\sigma(m)}{t}z^t\right)
=(1-z)^{-\sigma(m)}.                                   \tag{12}\label{eq:zeta}
\]
This is recorded for completeness; once the fixed-root census is known, it
is generic cycle bookkeeping and receives no standalone credit.

\section{Exact inverse signature}

Let
\[
S(n)=\bigl(a,\sigma(m),\tau(n)\bigr),\qquad n=2^am,\quad m\text{ odd}.
\]
If \(S(n)=S(q)\), match the distinguished roots, match all other roots
arbitrarily, preserve the binary positions in \eqref{eq:lifts}, and match the
\(\tau(n)\) extra rotation leaves.  This constructs a bijection commuting
with the two maps, so equal signatures are sufficient.

Conversely, the unlabelled graph recovers \(\sigma(m)\) as its number of
fixed vertices.  Let \(L\) be its maximum tail.  If \(L\geq2\), then
\(a=L\).  The only ambiguity is \(L=1\), where \(a\in\{0,1\}\).  Because
\(n\geq3\), in either case the relevant odd part exceeds one and there is
more than one fixed root.  If \(a=0\), only the distinguished root has
nonfixed predecessors.  If \(a=1\), every root has two level-one dihedral
predecessors.  Hence the graph distinguishes the two cases and recovers
\(a\).

Now the complete population of dihedral vertices is
\[
\sigma(m)(1+2+\cdots+2^a)
=\sigma(m)(2^{a+1}-1).
\]
The remaining vertices are exactly the rotations, so the graph recovers
\[
\tau(n)=|V|-\sigma(m)(2^{a+1}-1).                     \tag{13}\label{eq:taurecovery}
\]
This total-vertex formula is essential: a shortcut comparing the
distinguished root with other roots would not cover the one-root case
\(m=1\).  Equation \eqref{eq:taurecovery} covers it without exception and
completes necessity in \eqref{eq:signature}.

\paragraph{The \(33/35\) collision.}
Both parameters are odd and
\[
\sigma(33)=48=\sigma(35),\qquad \tau(33)=4=\tau(35).
\]
All 48 dihedral vertices are fixed.  Match \(H_{1,0}\) with \(H_{1,0}\),
match the remaining \((d,j)\) pairs lexicographically, and match the four
rotation leaves by
\[
R_1\mapsto R_1,\qquad R_3\mapsto R_5,\qquad
R_{11}\mapsto R_7,\qquad R_{33}\mapsto R_{35}.         \tag{14}\label{eq:collision}
\]
Every rotation maps to the distinguished root, so this 52-state bijection
commutes with the updates.  Multiplication by \(2^b\) preserves
\(a=b\), the odd-part divisor sum \(48\), and
\(\tau(2^b33)=4(b+1)=\tau(2^b35)\); \eqref{eq:signature} gives every lifted
collision.

\paragraph{Sharpness and boundary examples.}
The third signature coordinate cannot be dropped.  For example,
\[
v_2(15)=v_2(23)=0,\qquad \sigma(15)=\sigma(23)=24,
\]
but \(\tau(15)=4\) and \(\tau(23)=2\); the graphs have 28 and 26 vertices.
At the other extreme, \eqref{eq:collision} shows that even the complete graph
does not identify the ambient group order.  We use the literal subgroup set,
not conjugacy classes of subgroups: quotienting by conjugacy merges residue
positions and changes the binary fibres.  Finally, the restriction \(n\geq3\)
keeps the standard nondegenerate dihedral convention; no statement here
uses \(n=1,2\) to resolve the \(a=0/1\) case.

\section{Owner subtraction, collision firewall, and exact audit}

\begin{table}[H]
\centering
\caption{Owner subtraction and internal system firewall.}
\label{tab:firewall}
\small
\begin{tabular}{L{0.20\linewidth}L{0.34\linewidth}L{0.37\linewidth}}
\toprule
Source or neighbour & Material receiving no separation credit
& Literal separation retained here\\
\midrule
\citet{Cavior1975,ConradDihedralII}
& Complete subgroup coordinates and counts
& None until the coordinates are iterated as one global map\\
\citet{Frenkel2003,ShelashEtAl2023}
& Odd/even one-step normalizer and halving
& Full forest, all-time target fibres, iff signature, and order collisions\\
P91 / P135
& Finite-group or centralizer vocabulary
& No relation shift, spectrum, wreath product, or partition orbit; the
carrier here is every subgroup and the arrow is its ambient normalizer\\
P142 reserve
& Divisors and two-adic halving
& Literal subgroup normalizers and an unlabelled-graph inverse theorem, not a
divisor-only gcd map\\
P153
& Finite noninvertible branches
& Binary forests into fixed roots, not factorial arms into a translating
field cycle\\
\bottomrule
\end{tabular}
\end{table}

Related two-adic normalizer-length calculations, such as the Wielandt
subgroup setting of \citet{ShelashAshrafi2020}, are also treated as structural
adjacency, not as support for a contribution claim.  Formula searches used
dihedral-subgroup coordinates, normalizer towers, hypernormalizers, parity
halving, functional graphs, and the \(33/35\) arithmetic signature.  The
bounded search found the direct one-step owners above, but no source stating
the entire iterated graph--fibre--signature conjunction.  That bounded result
is not an ownership certificate.

\paragraph{Deterministic replay.}
The accompanying \texttt{verify.py} uses no randomness or external package.
For 44 values of \(n\), it materializes every displayed subgroup as an
element set, computes every normalizer by conjugating every subgroup element
with every ambient group element, and only then compares with
\eqref{eq:onestep}.  It next checks depths, fixed roots, images, and every
target fibre through three steps beyond stabilization.  Four equal-signature
pairs, beginning with \(33/35\), receive explicit commuting bijections.  The
frozen \texttt{CANONICAL.txt} ends with
\begin{verbatim}
PROFILE_SHA256 6eed12ce0c63f2d20f734ac1fa67634ce445140372dfc53e779a389de023b782
TOTAL boxes=44 iso_pairs=4 assertions=29590
VERDICT PASS_EXACT_REPLAY
\end{verbatim}
An independent pressure scan recorded no signature mismatch for
\(3\leq n\leq500\), but bounded enumeration is only a regression and
falsification control.  Subgroup completeness and both directions of
\eqref{eq:signature} are supplied by the proofs above.

\paragraph{Reproducibility and release state.}
The repository includes the exact script, canonical transcript,
claim--evidence matrix, source-verification ledger, controls, and deterministic
build instructions.  The manuscript is anonymous and uses no human-subject
data or randomized inference.  External circulation remains
\texttt{HOLD\_EXTERNAL}.

\bibliographystyle{plainnat}
\bibliography{references}

\end{document}
